% !PRINTED 251
\weberpage{251}
%<<CONT
those expressions we interchange $\alpha$ with $\beta$, then $\psi_1$, $\psi_2$, $\psi_3$ pass into $\psi_1$, $-\psi_3$, $-\psi_2$, and if we simultaneously interchange $\alpha$ with $\beta$ and $\gamma$ with $\delta$, then $\psi_1$, $\psi_2$, $\psi_3$ pass into $\psi_1$, $-\psi_2$, $-\psi_3$. From this the following results:

If one sets an arbitrary real value for the variable $x$, and if the roots $\alpha$, $\beta$, $\gamma$, $\delta$ of the biquadratic form are all real, then $\psi_1$, $\psi_2$, $\psi_3$ are also real.

If $\alpha$, $\beta$ are conjugate imaginary and $\gamma$, $\delta$ real, then the interchange of $\alpha$ and $\beta$ is equivalent to the interchange of $i$ with $-i$. In this case, therefore, $\psi_1$ is real, and $\psi_2$, $\psi_3$ are conjugate imaginary.

If, finally, $\alpha$, $\beta$ and $\gamma$, $\delta$ are two pairs of conjugate imaginary roots, then $\alpha$ is interchanged with $\beta$ and $\gamma$ with $\delta$ when $i$ passes into $-i$; hence in this case $\psi_1$ is real, and $\psi_2$, $\psi_3$ are \gesperrt{purely imaginary}.

If four roots are real, then $\psi_1^2$, $\psi_2^2$, $\psi_3^2$ become real and positive.

If two roots are real, then $\psi_1^2$ is real and positive, and $\psi_2^2$, $\psi_3^2$ are conjugate imaginary.

If four roots are imaginary, then $\psi_1^2$ is positive, and $\psi_2^2$, $\psi_3^2$ are real and negative.

Now in \S.\,66, (7) we have set up the cubic equation whose roots are $\psi_1^2$, $\psi_2^2$, $\psi_3^2$, and have also formed its discriminant. In the previous section we became acquainted with a criterion for the number of positive roots of a cubic equation, from which the following result is obtained:

If $D < 0$, the biquadratic equation has two real and two imaginary roots.

If $D > 0$, then, if four real roots are to exist, in the sequence
\[
1, \quad H, \quad H^2 - 16 A f^2, \quad -T^2
\]
three changes of sign must occur, i.e.\ we must have
\[
H < 0, \quad H^2 - 16 A f^2 > 0
\]
hold. In the three other still possible combinations of signs all four roots are imaginary. Here, however, an arbitrary real value may be set for $x$.
